\section{Introduction}
Software evolution histories provide a rich source of evidence for understanding how software systems change over time. Researchers use historical data to investigate refactoring practices~\cite{vassallo2019large}, developer collaboration patterns~\cite{kalliamvakou2015open}, technical debt accumulation and repayment~\cite{martini2015investigating,chowdhury_evidence_2025,digkas2018developers}, and the co-evolution of related software artifacts~\cite{zaidman2011studying,wang2021understanding}. Insights from these studies inform the development of maintenance models and tools, including automated refactoring~\cite{baqais2020automatic}, bug prediction~\cite{Wattanakriengkrai:2022,wang2022bugpre,PASCARELLA:2020,Rahman:2017}, change prediction~\cite{arvanitou2017method,chowdhury2022empirical,chowdhury:2022,Shihab:2012}, and automated technical debt repayment~\cite{mastropaolo2023towards}. Beyond research, practitioners rely on evolution histories for provenance analysis, traceability, developer onboarding, code comprehension, and expert identification~\cite{grund:2021}.

Because evolution histories are typically mined from version control systems, most prior work has examined change at coarse levels of granularity, including commits, files, and classes~\cite{negara2012dangerous}. However, researchers and practitioners increasingly require more fine-grained histories~\cite{Ryosuke:2019,Li:2018,grund:2021}. Class-level maintenance models, for example, are often less actionable because locating specific issues within a large class or file is difficult~\cite{PASCARELLA:2020,grund:2021,Mo:2022}. While line-level history reconstruction has been explored~\cite{Canfora:2007,Chen:2001}, its accuracy is limited by spurious matches between syntactically similar lines~\cite{Daniela:2014,Francisco:2017}. To address these limitations, recent research has focused on method-level history reconstruction. Tools such as \emph{FinerGit} (JSS 2020)~\cite{HIGO:2020}, \emph{CodeShovel} (ICSE 2021)~\cite{grund:2021}, \emph{CodeTracker} (FSE 2022)~\cite{jodavi_accurate_2022}, and more recently \emph{HistoryFinder} (ICSE 2026)~\cite{islam_historyfinder_2025} can trace methods across renames, moves, splits, and other transformations, enabling more precise studies of change patterns, defect proneness, and technical debt survival at the method level~\cite{chowdhury_evidence_2025,chowdhury2024method,chowdhury2025good}.

These history tracking tools, however, were developed and evaluated exclusively on production methods. Consequently, test evolution research remains constrained by the lack of reliable test method histories. Pinto \textit{et al.}~\cite{pinto2012understanding}, for example, reported challenges recovering test method histories because test methods are frequently renamed or moved. Bhatta \textit{et al.}~\cite{bhatta2025understanding} combined custom heuristics with \emph{RefactoringMiner}~\cite{tsantalis2020refactoringminer} and extensive manual inspection to identify deleted test methods, limiting scalability. Similarly, Spadini \textit{et al.}~\cite{Spadini:2018} relied on an unevaluated heuristic-based history construction approach when studying the relationship between test evolution and smells. Together, these studies suggest that the absence of a validated test method history tracking remains a key obstacle for understanding test code evolution.

Reliable test method histories, however, represent only part of the solution. Understanding \emph{why} test methods evolve also requires identifying the production methods they exercise. Without such mappings, it is difficult to determine whether a test changed in response to production code evolution or due to factors internal to the test itself, such as poor design or test smells. Method-level production-to-test mappings would also enable more precise studies of production--test co-evolution, which have largely been conducted at file or class granularity~\cite{zaidman2011studying,sun_revisiting_2023,wang2021understanding}. Furthermore, they could support finer-grained detection of obsolete tests by identifying the specific test methods affected by changes to particular production methods~\cite{bhatta2025understanding,yaraghi2025automated,pinto2012understanding,hao2013bug}.

To address these methodological challenges, we first answer the following two research questions. 
% first evaluate existing production code-based approaches for test method history tracking and production-to-test mapping. 

% We then leverage the most effective approaches to conduct the first large-scale study of test code evolution at the method level. Specifically, we investigate four research questions.

\emph{RQ1: Can existing history tracking tools effectively track test method revision histories?} We manually constructed an oracle of 120 test-method revision histories from 40 popular open-source Java projects. We then evaluate state-of-the-art history tracking tools originally developed for production code. The results show that these tools generalize well to test methods without recalibration. 

\emph{RQ2: What is the most accurate approach for method-level production-to-test mapping?} In addition to existing oracles, we manually constructed a new oracle from 20 projects and used the combined benchmark to evaluate existing mapping approaches. Our results show that no single approach dominates: some achieve high precision but low recall, while others favor high recall at the expense of precision.

Building on the findings from \textbf{RQ1} and \textbf{RQ2}, we revisit test code evolution and its relationship with test smells through the following two research questions.

\emph{RQ3: Are test methods revised less frequently?} Using the most effective history tracking and mapping approaches, we compare the revision frequencies of test methods and the production methods they exercise. Contrary to the common assumption, we find that many test methods undergo more revisions than their corresponding production methods.

\emph{RQ4: Are test smells associated with high revisions?} We investigate whether test smells are more associated with test methods that accumulate more revisions than the production methods they exercise, helping explain maintenance effort beyond what can be attributed to production-test co-evolution. Contrary to earlier findings~\cite{Spadini:2018}, we show that smells in small methods are more harmful than smells in large methods. 


In summary, our study makes three primary contributions. First, we provide the first systematic evaluation of state-of-the-art method history tracking tools on test methods. Second, we evaluate the existing method-level production-to-test mapping techniques using the existing and a new manually constructed oracle. Third, we conduct the first large-scale method-level study enabled by these validated techniques and provide new empirical insights into how test methods evolve relative to production methods and the role of test smells in test evolution. 

\textbf{Replication Package.} We share all data, oracles, and source code to enable replication and future research.\footnote{\href{https://github.com/SQMLab/method-co-evolution/blob/main/replication-package.md}{https://github.com/SQMLab/method-co-evolution/blob/main/replication-package.md}} This package also includes our manually annotated labels, which required approximately 180 person-hours to construct.

% Software evolution histories provide a rich source of evidence for understanding how software systems change over time. Researchers use historical data to investigate refactoring practices~\cite{vassallo2019large}, developer collaboration patterns~\cite{kalliamvakou2015open}, technical debt accumulation and repayment~\cite{martini2015investigating,chowdhury_evidence_2025,digkas2018developers}, and the co-evolution of related software artifacts~\cite{zaidman2011studying,wang2021understanding}. These insights inform the development of effective maintenance models and tools, including automated refactoring~\cite{baqais2020automatic}, bug prediction~\cite{Wattanakriengkrai:2022, wang2022bugpre, PASCARELLA:2020, Rahman:2017}, change prediction~\cite{arvanitou2017method,chowdhury2022empirical,chowdhury:2022, Shihab:2012}, and automated technical debt repayment~\cite{mastropaolo2023towards}. Beyond research, practitioners rely on evolution histories for provenance analysis, traceability, developer onboarding, code comprehension, and expert identification~\cite{grund:2021}.

% Because evolution histories are typically mined from version control systems, most prior work has examined change at coarse levels of granularity, including commits, files, and classes~\cite{negara2012dangerous}. However, researchers and practitioners often require more granular change histories~\cite{Ryosuke:2019,Li:2018,grund:2021}. Class-level maintenance models, for example, are less actionable because locating specific issues within a large class or file remains difficult~\cite{PASCARELLA:2020,grund:2021,Mo:2022}. While line-level history reconstruction has been attempted~\cite{Canfora:2007, Chen:2001}, its accuracy is limited by spurious matches between syntactically similar lines~\cite{Daniela:2014,Francisco:2017}. To address this limitation, recent work has focused on method-level history reconstruction. Tools such as \emph{FinerGit} (JSS 2020)~\cite{HIGO:2020}, \emph{CodeShovel} (ICSE 2021)~\cite{grund:2021}, \emph{CodeTracker} (FSE 2022)~\cite{jodavi_accurate_2022}, and, more recently, \emph{HistoryFinder} (ICSE 2026)~\cite{islam_historyfinder_2025} can trace methods across renames, moves, splits, and other transformations, enabling more precise studies of change patterns, defect proneness, and technical debt survival at the method level~\cite{chowdhury_evidence_2025,chowdhury2024method,chowdhury2025good}.

% Unfortunately, these method-level history tracking tools were trained and evaluated exclusively on production code. Consequently, the lack of reliable test method histories has limited the precision and scalability of test evolution research on test code evolution, despite its critical role in software maintenance. Pinto \textit{et al.}\cite{pinto2012understanding}, for instance, reported significant challenges recovering test method histories because test methods are frequently renamed or moved. Bhatta \textit{et al.}\cite{bhatta2025understanding} combined custom heuristics, RefactoringMiner~\cite{tsantalis2020refactoringminer}, and extensive manual inspection to identify deleted test methods, limiting the accuracy and scale of their analysis. Similarly, Spadini \textit{et al.}~\cite{Spadini:2018} relied on an unevaluated heuristic-based approach to reconstruct test method histories when studying the relationship between test evolution and test smells.

% Tracking test method evolution alone, however, is insufficient to understand \emph{why} test methods change. We also need to link test methods to the production methods they exercise. Without such mappings, it is difficult to determine whether a test changed in response to production code evolution, or other factors such as test smells. Production-to-test mappings at the method level would also advance research on production–test co-evolution, which has largely been conducted at coarser granularities~\cite{zaidman2011studying,sun_revisiting_2023,wang2021understanding}. Additionally, this mapping would be useful for more granular detection of outdated or obsolete tests~\cite{bhatta2025understanding} requiring test repairs due to changes in the production code by pinpointing which specific test methods in a test class became obsolete after some particular production methods changed in a production class~\cite{yaraghi2025automated,pinto2012understanding,hao2013bug}.

% Based on the importance of accurate change tracking and production-to-test mapping at the method-level granularity, we investigate which existing history tracking tools perform best for test methods, and which production-to-test mapping approach offers the best accuracy. We then apply these approaches to how test methods evolve when compared to source methods, and whether test smells are associated with test methods that change even more frequently than their corresponding source methods. Our contributions are organized around four research questions.

% \emph{RQ1: Can existing history tracking tools effectively track test method revision histories?} We manually constructed an oracle containing the revision histories of 120 test methods from 40 popular open-source Java projects. Using this oracle, we evaluated state-of-the-art history tracking tools that were trained and tested on production code only and found that all achieve high accuracy on test methods. However, \emph{HistoryFinder}~\cite{islam_historyfinder_2025} provides the best combination of accuracy and efficiency, making it our tool of choice for subsequent analyses.

% \emph{RQ2: What is the most accurate approach for method-level production-to-test mapping?} In addition to using the existing oracles, we manually constructed a new oracle of XXX production-to-test method mappings. We then evaluated the existing mapping techniques and found that there is no silver bullet: some tools achieve high precision but low recall, while others offer the opposite.

% \emph{RQ3: Are test methods revised less frequently than production methods?} Contrary to the prevailing view that test code is maintained less frequently, lags behind production changes, and generally evolves in response to changes in their production code, we found that approximately 30\% of test methods underwent more revisions than the corresponding production methods they exercise. This finding, enabled by method-level history tracking and mapping, reveals a level of test maintenance activity that has remained hidden at coarser granularities.

% \emph{RQ4: Are test smells associated with excessive test method revisions?} We found that test smells are more associated with the group of test methods that accumulate disproportionately more revisions than their production counterparts. This suggests that smells may be an important driver of test maintenance effort and reinforces the need for effective detection and remediation tools.

% \textbf{Replication Package.} \textcolor{red}{We share all the data and code to enable replication and extension.} 
% \section{Introduction}
% Software evolution histories provide a rich source of evidence for understanding how software systems change over time. Researchers use historical data to understand refactoring practices~\cite{vassallo2019large}, developer collaboration patterns~\cite{kalliamvakou2015open}, technical debt's accumulation, impact, and repayment~\cite{martini2015investigating,chowdhury_evidence_2025, digkas2018developers}, the co-evolution of related software artifacts~\cite{zaidman2011studying, wang2021understanding}. These understandings are important for developing effecting maintenace models and tools, including automated refactoring, bug prediction, change prediction, and automated technical debt repayment. Beyond research, practitioners rely on evolution histories for provenance analysis, traceability, developer onboarding, code comprehension, and expert identification~\cite{grund:2021}. 

% Because evolution histories are typically mined from version control systems, most prior work has examined change at coarse levels of granularity, such as commits, files, classes, and features. However, researchers and practitioners are often interested in more granular change history. For example, practitioners find class-level maintenance models less useful due to the difficulty in locating maintenance issues in a large class or file~\cite{PASCARELLA:2020, grund:2021, Mo:2022}. While some researchers aimed to generate line-level change history, their accuracy is limited because multiple lines can be similar just by chance.

% To address this problem, recent work has focused on method-level history reconstruction. Tools such as \emph{FinerGit} (JSS 2020), \emph{CodeShovel} (ICSE 2021), \emph{CodeTracker} (FSE 2022), and more recently, \emph{HistoryFinder} (ICSE 2026) can trace methods across renames, moves, splits, and other transformations. These advances enable more precise studies of method change patterns, defect proneness, and survival analysis of technical debt~\cite{chowdhury_evidence_2025, chowdhury2024method, chowdhury2025good}. Unfortunately, all these method-level history generation tools were trained and tested on production code only.

% Similar to production code, understanding test code evolution is instrumental for effective software maintenance. We argue that the lack of means to track change history for test methods has constrained our understanding of test code evolution. For example, in their attempt to understand the myths and realities of test suite evolution, Pinto \textit{et al.} had to recover change history of test methods\cite{pinto2012understanding}, but found it challenging because test methods are often renamed or moved. Similarly, while studying the reasons and impact of test deletion in Java, Bhatta \textit{et al.}\cite{bhatta2025understanding} had to combine custom heuristics, RefactoringMiner, and laborious manual inspection to identify deleted test methods, limiting the number of samples they could have analyzed. Similarly, Spadini \textit{et al.}~\cite{Spadini:2018} relied on an unevaluated heuristic-based history reconstruction approach to study test methods evolution with test smells. 

% We also argue that tracking test method evolution alone, however, is not enough to understand why a test method evolves. We need to link test methods to the production methods they exercise. Without such production-to-test mappings, it is difficult to determine whether a test changed in response to production code evolution, independent maintenance activities, or other factors, such as test smells. This linking is also important to improve our understanding of production--test co-evolution, as previous research in this line focused on higher granularity than method level.~\cite{zaidman2011studying,sun_revisiting_2023,wang2021understanding}. This linking can also help us improve test obsolescence research, because existing studies usually work at the file or class level, aiming to identify the test class that becomes obsolete due to some changes in its production class. Maintenance at such a high-level granularity is often less useful for practical applications, because it cannot reveal which specific test method became obsolete after a change to a particular production method, leaving developers to identify the affected test manually from the affected class, which is often very large to work with~\cite{grund:2021,chowdhury2024method, Mo:2022}. 

% Motivated by these two important observations, we investigate whether existing production code-based method history trackers can be used directly for tracking test method change history (RQ1), and which existing mapping approaches offer the best accuracy while mapping test methods to their corresponding production method (RQ2). Using our findings, we then proceed to accurately understand if test methods evolve primarily by the changes in their source methods, and if not, can test smells, to some extent, explain the undesired test method changes. In general, our contributions can be summarized by the following four research questions. 

% \emph{RQ1: Can existing history tracking tools for production methods be effectively used to track revision histories of test methods?}
% To answer this question, we manually constructed an oracle containing the revision histories of 120 test methods from 40 popular open-source Java projects. Using this oracle, we evaluate state-of-the-art method history tracking tools and assess their applicability to test code. Our results show that all evaluated tools achieve high accuracy on test methods, with \emph{HistoryFinder} providing the best combination of accuracy and efficiency.

% \emph{RQ2: What is the most accurate approach for production-to-test mapping at the method level?}
% We manually constructed a new oracle of production-to-test method mappings and used it to evaluate existing mapping approaches. Our results reveal that no single technique dominates across all evaluation criteria: some approaches achieve high precision, while others provide substantially higher recall.

% \emph{RQ3: Are test methods revised less frequently than production methods?}
% We observed many projects in which test methods exhibited greater evolutionary activity than production methods. More surprisingly, pairwise comparisons revealed that approximately 30\% of test methods underwent substantially more revisions than the corresponding production methods they exercise. This evidence, enabled by method-level history tracking and mapping, challenges the prevailing view in the co-evolution literature that test code is maintained less frequently and typically lags behind production code changes.

% \emph{RQ4: Can test smells explain why some test methods are revised more often than their production methods?}
% A test method is generally expected to evolve in response to changes in the production method it exercises. We therefore investigated whether test smells are associated with cases where test methods accumulate substantially more revisions than their production counterparts. Our results show that highly change-prone test methods contain significantly more test smells, suggesting that test smells may be an important driver of test maintenance effort and reinforcing the need for effective smell detection and remediation techniques.

% Previous test evolution studies have focused on understanding test suite evolution~\cite{pinto2012understanding}, production--test co-evolution~\cite{zaidman2011studying,sun_revisiting_2023,wang2021understanding}, smells and their impact on test method revisions~\cite{Spadini:2018}, and test deletion or obsolescence~\cite{bhatta2025understanding,hu2023identify,hao2013bug}. 

% These studies were constrained by two important factors: i) the lack of reliable method-level test histories, and ii) absence of mappings between test methods and the production methods they exercise. 
% % Consequently, many fundamental questions about how individual test methods evolve, and how their evolution relates to the production methods they exercise, remain unanswered.  

% \emph{Problem 1: change history at the test method level.} 
% \emph{Problem 1: production-to-test mapping at the method level}.  Spadini \textit{et al.}~\cite{Spadini:2018} blamed test smells for highly change prone method without verifying if those methods were changed mainly as a response to the change in the production method they exercise. Also, in obsolete test identification research, 


% Tracking revision history alone is not enough to accurately understand why test methods evolve. We need to map test methods to their corresponding source methods, which will enable us to observe if test methods are evolving due to the change in their source methods (normal), or they are evolving on their own (abnormal). Also, 
% these studies were significantly constrained due to the studies have either been conducted at coarse granularities (e.g., commits, files, and classes), or at the method-level granuarlity with unreliable history data. For example, Pinto \textit{et al.}~\cite{pinto2012understanding} reported difficulties in tracking test methods across revisions due to frequent transformations such as method rename and move, while Bhatta \textit{et al.}~\cite{bhatta2025understanding} combined custom heuristics, RefactoringMiner, and laborious manual inspection to identify deleted test methods. 



% Two long-standing challenges have hindered method-level studies of test evolution. First, accurately reconstructing the revision histories of test methods is difficult because methods may be renamed, moved, split, merged, or otherwise transformed over time. Existing studies often rely on heuristics, refactoring detection tools, and substantial manual validation to recover test histories. Second, 


% Recent advances in method-level history reconstruction provide an opportunity to overcome the first challenge. Tools such as \emph{FinerGit} (JSS 2020), \emph{CodeShovel} (ICSE 2021), \emph{CodeTracker} (FSE 2022), and more recently \emph{HistoryFinder} (ICSE 2026) can trace methods across complex transformations including renames, moves, and splits. Yet it remains unclear whether these tools are equally effective on test methods, which often exhibit naming conventions, structures, and evolution patterns that differ from production code. Moreover, no comprehensive study has combined accurate method-level histories with precise production-to-test mappings to investigate test evolution at scale.

% In this paper, we argue that reliable method-level test histories and accurate production-to-test mappings enable a new generation of studies on test maintenance and production--test co-evolution. Using state-of-the-art history tracking techniques and method-level source-to-test mappings, we revisit several long-standing assumptions about test evolution. Most notably, we find that test methods are not always the relatively passive artifacts portrayed in earlier co-evolution studies. Across many projects, approximately 30\% of test methods underwent substantially more revisions than the corresponding production methods they exercise, revealing a level of evolutionary activity that has largely remained hidden at coarser levels of analysis.


% Software evolution histories are a valuable source of evidence for understanding how software systems change over time. Researchers use them to study refactoring, developer collaboration, technical debt, and the co-evolution of related artifacts, while practitioners rely on them for provenance analysis, traceability, onboarding, code comprehension, and expert identification~\cite{grund:2021}.

% Because these histories are typically mined from version control systems, most prior work has examined change at coarse levels of granularity, such as commits, files, classes, and features. However, many important questions require finer-grained histories. Line-level reconstruction approaches exist, but their accuracy is limited because syntactically similar lines can be mismatched across revisions. To address this problem, recent work has focused on method-level history reconstruction. Tools such as \emph{FinerGit} (JSS 2020), \emph{CodeShovel} (ICSE 2021), \emph{CodeTracker} (FSE 2022), and more recently, \emph{HistoryFinder} (ICSE 2026) can trace methods across renames, moves, splits, and other transformations. These advances have enabled more precise studies of method change patterns, defect proneness, refactoring behavior, and self-admitted technical debt.

% In contrast, test code evolution remains relatively underexplored, even though it is central to automated test repair, test maintenance, test smell management, regression testing, and obsolete test detection. Prior studies have examined test suite evolution~\cite{pinto2012understanding}, production--test co-evolution~\cite{zaidman2011studying,sun_revisiting_2023,wang2021understanding}, test smells~\cite{Spadini:2018,garousi2018smells,tufano2016empirical}, and test deletion or obsolescence~\cite{bhatta2025understanding,hu2023identify,hao2013bug}. Unfortunately, findings of these studies are constrained by two different factors: the lack of i) reliable method-level test histories, and ii) precise mappings between test methods and the production methods they exercise.

% \emph{Test method History.} Researchers often reconstruct method-level test histories using heuristics, refactoring detection tools, and substantial manual validation. Pinto \textit{et al.}\cite{pinto2012understanding} reported difficulties tracking test cases because test methods are often renamed or moved. Bhatta \textit{et al.}\cite{bhatta2025understanding} combined RefactoringMiner, custom heuristics, and manual inspection to identify deleted test methods. Spadini \textit{et al.}~\cite{Spadini:2018} relied on unevaluated heuristic-based history reconstruction to study test methods evolution. 

% \emph{Production-to-test mapping}. Tracking revision history alone is not enough to accurately understand why test methods evolve. We need to map test methods to their corresponding source methods, which will enable us to observe if test methods are evolving due to the change in their source methods (normal), or they are evolving on their own (abnormal). Also, existing test obsolescence studies usually work at the file or class level, which is often less useful for practical applications, because they cannot reveal which specific test method became obsolete after a change to a particular production method, leaving developers to identify the affected test manually.

% We argue that accurate method-level test histories in conjunction with precise method-level source-to-test mappings would enable more rigorous and practically useful new studies of test maintenance and production--test co-evolution, while also allowing earlier findings to be revisited with more reliable data. As such, this paper answers the following four research questions. 

% \emph{RQ1: Can existing history tracking tools for production methods be effectively used to track revision history of test methods?} We manually constructed the revision histories of 120 test methods originating from 40
% popular open-source Java projects, \textcolor{red}{requiring x person-hours in total.} With this oracle, we evaluated all the state-of-the-art history tracking tools that were trained and tested on production methods only.  Encouragingly, our evaluation suggests that all these tools exhibit very good accuracy on test methods as well. However, \emph{HistoryFinder} outperforms all other tools in both F1-score and execution time, making it an obvious choice for our subsequent analysis. 

% \emph{RQ2: What is the most accurate approach for production-to-test mapping at the method level?} To accurately identify the most accurate approach, we made two contributions. In addition to having the existing oracles of production-to-test examples, we manually created another new oracle of production-to-test methods, \textcolor{red}{requiring x person-hours}. We then evaluated the existing mapping approaches and found that there is no silver bullet. Some approaches show high precision, while others offer high recall.  

% \emph{RQ3: Are test methods revised less than production methods?} We observed many projects in which test methods exhibited greater evolutionary activity than production methods. More surprisingly, pairwise comparisons showed that approximately 30\% of test methods underwent substantially more revisions than the corresponding production methods they exercise. This evidence, enabled by method-level history tracking and mapping, challenges the prevailing view in the co-evolution literature that test code is maintained less frequently and typically lags behind production code changes.

% \emph{RQ3: Can test smell explain the many more revisions of a test method than its production method?}
% Normally, a test method should only require revision when its corresponding production method changes—so test revisions should not exceed production revisions. We investigated whether test smells explain cases where this expectation is violated. The answer is yes: test methods revised more frequently than their production counterparts contained significantly more smells, underscoring the need for better detection and remediation tools.
% \textbf{RQ4: Which test smells present when a test method is introduced are associated with its subsequent recurrent revision-proneness?}





% Software evolution histories provide a rich source of information for understanding how software systems change over time. Researchers leverage historical data to investigate phenomena such as refactoring practices, developer collaboration patterns, technical debt accumulation and repayment, and the co-evolution of related software artifacts. Insights from these studies have informed the development of effective software maintenance techniques and tools. Beyond research, practitioners use evolution histories for provenance analysis, traceability, developer onboarding, code comprehension, and expert identification~\cite{grund:2021}.

% Because software histories are typically extracted from version control systems, they have traditionally been studied at relatively coarse granularities, such as commits, files, classes, and features. However, researchers and practitioners often require evolution of finer-grained code components. Although several approaches have attempted to reconstruct line-level histories, their effectiveness is often limited because syntactically similar lines can be incorrectly matched across revisions. To overcome these limitations, recent research has focused on reconstructing method-level evolution histories. Tools such as \emph{FinerGit} (JSS 2020), \emph{CodeShovel} (ICSE 2021), \emph{CodeTracker} (FSE 2022), and, more recently, \emph{HistoryFinder} (ICSE 2026) can accurately trace methods across renames, moves, splits, and other transformations. By recovering complete method histories, these tools have enabled fine-grained studies of software evolution, including investigations of method change patterns, defect proneness, refactoring behavior, and the lifespan of self-admitted technical debt.

% While these advances have significantly improved our understanding of production code evolution, test code evolution remains comparatively underexplored. Yet, understanding how tests evolve is critical for a wide range of software engineering activities, including automated test repair, test maintenance, test smell management, regression testing, and obsolete test detection. Existing studies have examined test evolution from several perspectives, including the evolution of test suites~\cite{pinto2012understanding}, production–test co-evolution~\cite{zaidman2011studying,sun_revisiting_2023,wang2021understanding}, the impact of test smells on maintainability~\cite{Spadini:2018,garousi2018smells,tufano2016empirical}, and the consequences of test deletion and obsolescence~\cite{bhatta2025understanding,hu2023identify,hao2013bug}. These studies, however, are constrained by the lack of reliable method-level test evolution histories and accurate mappings between test and production methods.

% Due to the lack of a specialized test method history tracker, researchers often rely on heuristics, refactoring detection tools, and extensive manual validation to reconstruct test histories. Pinto \textit{et al.}~\cite{pinto2012understanding}, for example, reported difficulties tracking test cases because test methods are frequently renamed or moved. Bhatta \textit{et al.}~\cite{bhatta2025understanding} combined RefactoringMiner, custom heuristics, and laborious manual inspection to identify deleted test methods, while Spadini \textit{et al.}~\cite{Spadini:2018} relied on unverified heuristic-based history reconstruction when studying the evolution of test smells. %Consequently, the accuracy and completeness of the resulting histories remain uncertain, potentially affecting the validity of the conclusions drawn from them.

% The lack of method-level source-to-test mappings introduces a second challenge. Existing co-evolution studies typically operate at the file or class level because establishing relationships between production and test methods is considerably more difficult than mapping classes. However, class-level co-evolution provides only a coarse approximation of the underlying relationships. For example, it does not tell which specific test method became obsolete for a change in a production method, so practitioners have to manually find that specific test method from the test class that needs updating.

% We therefore argue that accurate test method histories, combined with precise source-to-test mappings at the method level, are necessary for advancing the study of test evolution. Such capabilities would not only enable new research directions in test maintenance and production–test co-evolution, but would also allow prior findings to be revisited using substantially more precise and reliable data. Motivated by this gap, this paper makes four contributions.

% Software evolution histories provide a rich source of information for understanding how software systems change over time. Researchers leverage historical data to investigate phenomena such as refactoring practices, developer collaboration patterns, technical debt accumulation and repayment, and the co-evolution of related software artifacts. Insights from these studies are important for building effective maintenance techniques and tools. Beyond research, practitioners use evolution histories for provenance analysis, traceability, developer onboarding, code comprehension, and expert identification~\cite{grund:2021}.

% Due to their dependence on version control systems, software evolution histories have traditionally been collected and analyzed at relatively coarse granularities, such as commits, files, classes, or features. However, researchers and practitioners often require understanding evolution at lower granularities. Although several approaches have attempted to reconstruct line-level histories, their effectiveness is often limited because syntactically similar lines can be matched incorrectly across revisions. To address this limitation, recent research has focused on reconstructing method-level evolution histories. Tools such as \emph{FinerGit} (JSS 2020), \emph{CodeShovel} (ICSE 2021), \emph{CodeTracker} (FSE 2022), and more recently \emph{HistoryFinder} (ICSE 2026) can accurately trace individual methods across renames, moves, splits, and other code transformations. By recovering complete method histories, these tools enable accurate fine-grained software evolution studies, including investigations of method change patterns, defect proneness, refactoring behavior, and the lifespan of self-admitted technical debt.

% Although the aforementioned previous studies predominantly focused on production code, understanding how tests evolve is also essential for many software engineering activities, including automated test repair, test maintenance, test smell management, regression testing, and obsolete test detection. Existing studies have examined test evolution from several perspectives. Researchers have investigated how and why test suites evolve~\cite{pinto2012understanding}, how production and test code co-evolve~\cite{zaidman2011studying,sun_revisiting_2023,wang2021understanding}, the impact of test smells on maintainability~\cite{Spadini:2018,garousi2018smells,tufano2016empirical}, and the consequences of test deletion~\cite{bhatta2025understanding} and obsolescence~\cite{hu2023identify,hao2013bug}. However, much of the existing literature on test code evolution is constrained by two important issues: i) lack of method-level test evolution data, and ii)  the absence of accurate mappings between production and test methods. 

% Due to the lack of a specialized test method history tracker, researchers often rely on a combination of heuristics, refactoring detection tools, and extensive manual validation to reconstruct test method histories. Pinto \textit{et al.}~\cite{pinto2012understanding}, for example, reported difficulties tracking test cases because test methods are frequently renamed or moved. Bhatta \textit{et al.}~\cite{bhatta2025understanding} combined RefactoringMiner, custom heuristics, and laborious manual inspection to identify deleted test methods, while Spadini \textit{et al.}~\cite{Spadini:2018} relied on unverified heuristic-based history reconstruction to study the evolution of test methods with and without smells. Likewise, while studying the co-evolution of test and production code, studies have focused on file/class level only, because mapping between test and production methods is more challenging than mapping between test and production classes---in a software project, there are many more methods than classes or files. Also, these studies assumed that if a test and its corresponding production class are modified in the same commit or the modifications happen within a small time difference, that change is a co-evolution. However, such assumptions incur many false positives. 

% We, therefore, argue that having a tool for accurate test method history generation and mapping between test to source at the method level would not only enable new research directions in test code maintenance, but also will help revisit the prior test studies with much more accuracy and confidence. Based on this, we make four concrete contributions in this paper. 



% We argue that 
% are driven by the assumption that test changes are primarily driven by changes in the production code they exercise, such assumption cannot be well verified without the absence of accurate history tracking of 
% this relationship is rarely examined at the method level. For example, while prior work has shown that smelly tests tend to undergo more modifications~\cite{Spadini:2018}, it remains unclear whether these changes originate from the tests themselves or are simply reactions to modifications in the corresponding production methods. More broadly, without method-level source-to-test mappings, it is difficult to distinguish intrinsic test evolution from evolution induced by production code changes.

% Combining accurate method-level test histories with precise source-to-test mappings would enable a fundamentally deeper understanding of test evolution. Such capabilities would support fine-grained studies of production–test co-evolution, facilitate root-cause analysis of test changes, improve obsolete test detection, and provide the foundation for a new generation of testing and maintenance tools.

% Researchers leverage code evolution histories to investigate a wide range of phenomena, including refactoring practices, developer collaboration patterns, technical debt accumulation and repayment, and the co-evolution of related software artifacts. Insights from these studies can enable the development of more effective tools and techniques for software maintenance. Beyond research, practitioners can use code histories for provenance analysis, traceability, developer onboarding, code comprehension, and expert identification~\cite{grund:2021}. 
% Unfortunately, code evolution data have predominantly been collected and analyzed at coarse-grained levels, such as features, commits, files, and classes, even though many applications require understanding software evolution at finer granularity that traditional version control systems cannot readily provide. While several approaches have attempted to reconstruct line-level histories, their precision and recall remain limited because similar lines of code can often be matched incorrectly by chance.

% To address these limitations, recent research has focused on developing techniques and tools for method-level history tracking, including \emph{FinerGit} (JSS 2020), \emph{CodeShovel} (ICSE 2021), \emph{CodeTracker} (FSE 2022), and, more recently, \emph{HistoryFinder} (ICSE 2026). By reconstructing the complete evolution history of individual methods, these tools enable researchers and practitioners to study software evolution research with unprecedented accuracy. For example, they facilitate accurate investigations into why methods evolve, how change-prone or defect-prone methods can be identified proactively, and how long self-admitted technical debt survives before being addressed. 

% Although test code is essential for ensuring software quality and facilitating maintenance activities, its evolution has received substantially less attention than that of production code. Nevertheless, understanding how test code evolves is important because many software engineering tasks—including automated test repair, test maintenance, test smell management, and obsolete test detection—depend on accurate knowledge of test evolution.

% Prior research has investigated test evolution from several perspectives. One line of work has focused on understanding how and why test suites evolve. For example, Pinto \textit{et al.}~\cite{pinto2012understanding} analyzed the evolution of test suites and reported that test modifications are often complex, posing challenges for automated test repair techniques. Another research direction has examined the relationship between production and test code evolution, motivated by the assumption that test changes are frequently triggered by modifications in the corresponding production code~\cite{zaidman2011studying,sun_revisiting_2023,wang2021understanding}. Researchers have also studied quality issues in test code, such as test smells, and their impact on maintainability and evolution. For instance, Spadini \textit{et al.}~\cite{Spadini:2018} found that smelly test code tends to be more change-prone and bug-prone than cleaner tests~\cite{garousi2018smells,tufano2016empirical}. More recently, researchers have investigated test removal and obsolete test code. Bhatta \textit{et al.}~\cite{bhatta2025understanding} studied deleted tests and found that, despite many deleted tests appearing redundant, their removal often leads to reductions in coverage and mutation scores. Other studies have explored techniques for identifying and repairing obsolete tests after production code changes~\cite{hu2023identify,hao2013bug}.

% Despite their valuable contributions, much of the existing literature is constrained by two fundamental limitations. First, most studies lack reliable support for tracking the evolution of individual test methods across software histories. Consequently, researchers often rely on heuristics, refactoring detection tools, or extensive manual validation to reconstruct test histories. For example, Pinto \textit{et al.}~\cite{pinto2012understanding} reported difficulties tracking test cases because tests are frequently renamed or moved. Similarly, Bhatta \textit{et al.}~\cite{bhatta2025understanding} combined RefactoringMiner with heuristics and manual verification to identify deleted tests, while Spadini \textit{et al.}~\cite{Spadini:2018} relied on heuristic-based history reconstruction to analyze test evolution. These challenges stem largely from the absence of specialized tools capable of accurately tracking test methods across commits.

% Second, existing studies rarely leverage source-to-test mappings at the method level. As a result, many observed test changes cannot be directly linked to the evolution of the production methods they are intended to validate. For example, although Spadini \textit{et al.}~\cite{Spadini:2018} associated test smells with increased test modifications, it remains unclear how many of those modifications were actually induced by changes in the corresponding production code. More generally, the lack of fine-grained source-to-test mappings limits our ability to distinguish between test changes driven by production evolution and those caused by other factors. Combining accurate method-level histories with precise source-to-test mappings would enable a deeper understanding of test evolution and support applications such as fine-grained co-evolution analysis, root-cause analysis of test changes, and method-level obsolete test detection, which is substantially more precise than existing class-level approaches.

% Although test code plays a critical role in software quality and maintenance, it has received considerably less research attention than production code. While investigating why and how a test suite evolves, Pinto \textit{et al.}~\cite{pinto2012understanding} found that test modifications are often so complex that existing approaches are not adequate for automated test repair techniques. Bhatta \textit{et al.}~\cite{bhatta2025understanding} studied test deletion and observed that although many redundant tests are deleted, they surprisingly reduce coverage and mutation scores. Multiple studies focused on production and test code co-evolution~\cite{zaidman2011studying,sun_revisiting_2023,wang2021understanding}, due to the common assumption that test changes are mainly induced by changes in their corresponding production code. While studying test smells~\cite{garousi2018smells, Spadini:2018,tufano2016empirical} and evolution, Spadini \textit{et al.}~\cite{Spadini:2018} observed that test code with smells is generally more change- and bug-prone than less smelly tests. Another area of testing research is to identify and fix obsolete test code once the corresponding production code is changed~\cite{hu2023identify,hao2013bug}. 

% We found that many of these test code evolution-related studies have suffered due to two main reasons: i) The absence of any specialized tools for accurately tracking test evolution at the method level. For example, in their test suite evolution study, Pinto \textit{et al.}~\cite{pinto2012understanding} had to capture whether a test case was found in two subsequent test suites, which was difficult because, according to their observations, many tests were moved or renamed, which made tracking them difficult. Similarly, in their test deletion study, Bhatta \textit{et al.}~\cite{bhatta2025understanding}, in addition to using RefactoringMiner, relied on heuristics and significant manual verification to build a dataset of deleted test cases. Likewise, in their study, Spadini \textit{et al.}~\cite{Spadini:2018} had to collect change history of test methods with unverified heuristics due to the absent of any specialized tools. ii) Not using source-to-test mapping at the method level. For example, Spadini \textit{et al.}~\cite{Spadini:2018} blamed test smells for higher revisions of test methods, without verifying if those revisions were done in response to the change in the corresponding source method. Also, joining accurate test method history with accurate source-to-test mapping would enable obsolete test detection at the method level, which would be much more efficient than the current approaches which detect test classes that need to be modified. 
%The few research that have studied test code evolution (or coevolution) focused on higher change granularity.




% Unfortunately, this insightful understanding at the method-level granularity is available for source (or production) code only, as all the method-level change tracking tools were trained and tested on source code. 



% We investigate the following research questions:

% \textbf{RQ1: Can existing method-history generation tools effectively generate test-method revision histories?}

% \textbf{RQ2: How accurately do existing test-to-production method linking approaches perform against different ground-truth datasets?}

% \textbf{RQ3: How do test and production method revisions differ?}

% \textbf{RQ4: Which test smells present when a test method is introduced are associated with its subsequent recurrent revision-proneness?}
